

\tiny
\begingroup
\renewcommand{\arraystretch}{1.24}
\begin{longtable}{>{\RaggedRight\arraybackslash}p{1.25cm} >{\RaggedRight\arraybackslash}p{1.55cm} >{\RaggedRight\arraybackslash}p{1.75cm} >{\RaggedRight\arraybackslash}p{1.05cm} >{\RaggedRight\arraybackslash}p{1.55cm} >{\RaggedRight\arraybackslash}p{1.9cm} >{\RaggedRight\arraybackslash}p{1.9cm} >{\RaggedRight\arraybackslash}p{2cm} >{\RaggedRight\arraybackslash}p{2.8cm} >{\centering\arraybackslash}p{1.15cm}}
\caption{Tabla consolidada de iniciativas, proyectos e instrumentos digitales identificados en los anexos}\label{tab:consolidado-proyectos-digitales}\\
\toprule
\textbf{Tema} & \textbf{Proyecto} & \textbf{Entidad responsable} & \textbf{Periodo} & \textbf{Presupuesto} & \textbf{Alcance geográfico} & \textbf{Beneficiarios} & \textbf{Objetivos} & \textbf{Principales componentes} & \textbf{Fuentes} \\
\midrule
\endfirsthead
\caption[]{Tabla consolidada de iniciativas, proyectos e instrumentos digitales identificados en los anexos (continuación)}\\
\toprule
\textbf{Tema} & \textbf{Proyecto} & \textbf{Entidad responsable} & \textbf{Periodo} & \textbf{Presupuesto} & \textbf{Alcance geográfico} & \textbf{Beneficiarios} & \textbf{Objetivos} & \textbf{Principales componentes} & \textbf{Fuentes} \\
\midrule
\endhead
\midrule
\multicolumn{10}{r}{\emph{Continúa en la siguiente página}}\\
\endfoot
\bottomrule
\endlastfoot
Redes móviles & Subasta 5G y obligaciones de hacer & MinTIC, ANE y operadores & 20-dic-2023 a 2026 & Recaudo de 1,37 billones de pesos; obligaciones de hacer por 389.711 millones de pesos & Nacional; obligaciones en instituciones educativas rurales y vías & Instituciones educativas rurales, usuarios de carreteras y usuarios móviles & Habilitar redes 5G y usar obligaciones de hacer para cerrar brechas de conectividad & Asignación de 320 MHz en banda de 3.500 MHz; bloques de 80 MHz; conectividad escolar y cobertura vial & \parencite{mintic_5g_realidad_2023} \\
Redes móviles & Inicio oficial del despliegue 5G & MinTIC y operadores & 23-feb-2024 & No especificado & Nacional & Usuarios móviles y operadores habilitados & Activar el despliegue de redes de quinta generación & Resoluciones que habilitan a operadores ganadores para desplegar redes 5G & \parencite{mintic_despliegue_5g_2024} \\
Redes móviles & Escuelas Potencia 5G & MinTIC y operadores & 2024--2026 & No especificado & 233 municipios en 21 departamentos & Hasta 1.191 instituciones educativas & Conectar instituciones educativas mediante obligaciones de hacer 5G & Aproximadamente 4.000 km de fibra y conectividad educativa de largo plazo & \parencite{mintic_escuelas_potencia_5g_2024} \\
Redes móviles & Controversia por adjudicatario 5G & MinTIC y Telecall & 29-ene-2025 & Más de 300.000 millones de pesos en discusión administrativa reportada & Nacional & Estado, operador adjudicatario y usuarios potenciales & Resolver obligaciones administrativas asociadas a la subasta 5G & Actuaciones por garantías y primer pago de contraprestación & \parencite{mintic_telecall_espectro_5g_2025,lorduy_telecall_conciliacion_2026} \\
Conectividad educativa & Centros Digitales & MinTIC / Fondo Único TIC & 2022--2033 & Contratos por 1,054 billones y 1,073 billones de pesos & 32 departamentos; corregimientos y veredas rurales & 14.057 comunidades e instituciones educativas, sobre todo rurales & Garantizar conectividad educativa gratuita y de largo plazo & Contratos 1042-2022 y 749-2022; operación de centros digitales; incremento progresivo de velocidad & \parencite{mintic_informe_gestion_molina_2025} \\
Conectividad PDET & Zonas Comunitarias para la Paz & MinTIC y ART & 2023--jul. 2026 & Convenio por 205.992 millones de pesos & 161 municipios PDET & 192.000 estudiantes y docentes; impacto indirecto en 12 millones de habitantes & Conectar sedes educativas y comunidades en zonas históricamente marginadas & 1.262 zonas WiFi instaladas; operación y mantenimiento; dos puntos de acceso por zona & \parencite{mintic_informe_gestion_molina_2025} \\
Conectividad hogares & ConectiVIDAd para Cambiar Vidas & MinTIC / Fondo Único TIC e Internexa & 2023--2026 & 270.588 millones de pesos para AE2; otros componentes separados & Antioquia-Urabá, Cauca, Chocó, La Guajira, Nariño y Valle del Cauca; nuevas convocatorias en Cauca, Chocó y Nariño & Hogares estratos 1 y 2, comunidades organizadas e instituciones educativas públicas & Ampliar Internet de banda ancha con red troncal y última milla en zonas desconectadas & Red troncal, más de 308 km de fibra óptica, expansión a hogares, WiFi y redes comunitarias & \parencite{mintic_informe_gestion_molina_2025} \\
Conectividad comunitaria & Juntas de Internet / Comunidades de Conectividad & MinTIC / Fondo Único TIC & 2024--2026 & 189.705 millones de pesos & 31 departamentos & 48.000 hogares aproximadamente; organizaciones comunitarias rurales & Convertir organizaciones locales en proveedoras de Internet fijo comunitario & 520 JI-CDC; infraestructura de red de acceso; operación, administración y sostenibilidad comunitaria & \parencite{mintic_juntas_internet_2024,mintic_informe_gestion_molina_2025} \\
Conectividad PDET / frontera & Catatumbo y Área Metropolitana de Cúcuta JI-CDC & MinTIC y convenio territorial 1211-2025 & 2025--2026 & 47.000 millones de pesos; 42.000 millones de aporte MinTIC & 17 municipios del Catatumbo y área metropolitana de Cúcuta & 6.000 hogares, 110 instituciones educativas y 900 beneficiarios de apropiación digital & Llevar conectividad satelital y comunitaria a una de las zonas más rezagadas y complejas del país & 100 comunidades de conectividad; servicio a hogares por 12 meses; servicio a escuelas por 18 meses; apropiación a juntas & \parencite{mintic_informe_gestion_molina_2025} \\
Conectividad rural comunitaria & Conectando a los no conectados & MinTIC, Unión Europea y Colnodo & 2025--2026 piloto & No especificado & Veredas Tasmag, El Guadual y Colinas Bajo, en Nariño y Guaviare & Comunidades rurales apartadas & Consolidar redes comunitarias de Internet para democratizar la tecnología & Nodos por red, centro de inclusión digital, computadores, servidor local y apropiación digital & \parencite{mintic_conectando_no_conectados_2025} \\
Acceso asequible & Líneas de Fomento 2.0 / 3.0 & MinTIC / Fondo Único TIC & 2024--2025 & 27.952 millones de pesos para 2.0; 3.0 no especificado & Nacional en 2.0; Ciudad Bolívar, Bogotá, en 3.0 & Hogares vulnerables e ISP regionales con menos de 30.000 usuarios & Ampliar acceso fijo residencial de banda ancha con operadores pequeños y focalización social & Convocatoria a PRST minoristas; interventoría; 4.500 hogares vulnerables focalizados en Ciudad Bolívar & \parencite{mintic_informe_gestion_molina_2025} \\
Territorial / OXI & Obras por Impuestos digitales & MinTIC, DNP, ART y contribuyentes privados & Aprobados 2025 & 49.318 millones de pesos en Nariño-Cauca; 6.439 millones de pesos en Tame & Nariño, Cauca y Tame (Arauca) & 180 escuelas, 157 entidades públicas, 194 zonas WiFi, juntas de Internet y más de 1.200 beneficiarios directos en Tame & Financiar infraestructuras de conectividad y apropiación en regiones apartadas vía OXI & Fibra troncal, nodos de distribución, última milla, WiFi, conectividad gubernamental y apropiación & \parencite{mintic_informe_gestion_molina_2025} \\
Radar & Red nacional de radares 3D & MinTransporte / Aerocivil & 2023--2026 & 332.717 millones de pesos & País completo, con sensores en Cali, Urabá, Leticia y despliegues previstos en San Andrés, Villavicencio, Riohacha y Bogotá & Aviación civil y militar; usuarios del espacio aéreo nacional & Renovar y modernizar la red de vigilancia aérea para mejorar navegación, seguridad operacional, supervisión y control del espacio aéreo & Estado: en ejecución; contrato Radcol 2023; red de siete radares 3D con radar primario 3D, secundario Modo S, ADS-B, canal meteorológico, torre de 30 m, cobertura entre 100/120 y 250 millas náuticas; Cali activado en julio de 2025, Urabá en septiembre de 2025, Leticia en enero de 2026; red prevista para finales de 2026 & \parencite{mintransporte_radar_3d_cali_2025,cepeda_radares_3d_colombia_2026,aerocivil_radar_uraba_2025} \\
Infraestructura educativa digital & Computadores para Educar / Tecnologías para el Aprendizaje & Computadores para Educar & 2025 & Superior a 170.000 millones de pesos & Nacional & Establecimientos educativos públicos; niños, niñas y jóvenes & Ampliar dotación tecnológica e innovación educativa & Convocatoria 2025; más de 14.400 equipos adquiridos; entregas directas y por sede; reutilización y robótica educativa & \parencite{computadores_para_educar_tecnologias_aprendizaje_2025} \\
IA territorial & Centros PotencIA & MinTIC y Findeter & Sep. 2024--abr. 2026 & 234.400 millones de pesos & 60 centros en 28 departamentos & Comunidades territoriales, jóvenes, emprendedores y entidades locales & Construir una red nacional de inclusión digital y apropiación de tecnologías emergentes e IA & Aulas de IA, laboratorios de prototipado, internet gratuito de alta velocidad, zonas colaborativas, XR, dotación y obra civil & \parencite{mintic_informe_gestion_molina_2025} \\
IA territorial & Centros PotencIA & MinTIC & 2024--2025 & 132.000 millones de pesos para diseños reportados en Usme y Zipaquirá & Meta nacional de 60 centros & Comunidades territoriales y usuarios de centros de IA & Ampliar apropiación territorial de IA y tecnologías emergentes & Diseños de centros de IA, conectividad, laboratorios y espacios de apropiación & \parencite{mintic_rendicion_cuentas_2024,mintic_centros_potencia_2025} \\
Formación avanzada & Talento Tech & MinTIC & 2024--jul. 2026 & 403.181 millones de pesos en Gobierno del Cambio; 161.566 millones de pesos en 2025 & 9 regiones & Meta mínima de 100.032 personas y meta de 113.070 certificadas & Reducir brecha de competencias digitales con bootcamps & Programación, IA, analítica, blockchain, nube y ciberseguridad; cuatro fases y 14 contratos & \parencite{mintic_talento_tech_lanzamiento_2023,mintic_talento_tech_regiones_2023,mintic_talento_tech_v2_2023, mintic_informe_gestion_molina_2025} \\
Formación masiva & SENATIC & MinTIC, SENA y OIT & Dic. 2023--dic. 2026 & 196.314 millones de pesos & Nacional & Meta de 303.900 personas certificadas & Fortalecer competencias digitales, empleabilidad y productividad & Certificación masiva; alianza con SENA y OIT; más de 144.000 personas certificadas a agosto de 2025 & \parencite{mintic_informe_gestion_molina_2025} \\
Niñez y juventud & Colombia Programa & MinTIC & Nov. 2023--jul. 2026 & 43.043 millones de pesos & Nacional & 11.200 docentes y 896.000 estudiantes de colegios oficiales & Llevar pensamiento computacional y programación desde preescolar a grado 11 & Centros de referencia; micro:bit; kits Biobots y STEM; mentorías; ferias Código en Acción; red de nodos & \parencite{mintic_informe_gestion_molina_2025} \\
Formación continua & AvanzaTec & MinTIC y empresas tecnológicas vía MoU & Jun. 2024--dic. 2026 & No especificado & Nacional & Meta de 97.000 personas certificadas & Ofrecer formación corta y certificaciones con aliados tecnológicos & Más de 45 ofertas; rutas básica, intermedia y avanzada; más de 45.000 certificaciones al corte 2025 & \parencite{mintic_informe_gestion_molina_2025} \\
Política digital & Estrategia Nacional Digital & DNP, Presidencia y MinTIC & 2023--2026 & No especificado & Nacional & Estado, ciudadanía y economía digital & Orientar la transformación digital pública, productiva y social & Conectividad, datos, seguridad, talento digital, IA, transformación pública y economía digital & \parencite{mintic_estrategia_nacional_digital_2023} \\
Política de IA & Hoja de Ruta de IA & MinCiencias & Febrero de 2024 & No especificado & Nacional & Ecosistema público, científico, productivo y social de IA & Definir un marco estratégico para adopción ética y sostenible de IA & Cinco entornos: ética y gobernanza; educación, investigación e innovación; industrias; datos; privacidad, ciberseguridad y defensa & \parencite{minciencias_hoja_ruta_ia_pdf_2024,minciencias_informe_gestion_2024_pdf} \\
Participación territorial en IA & Diálogos Regionales en IA & MinCiencias & Abril--mayo de 2024 & No especificado & Tumaco, Bucaramanga, Villavicencio, Medellín, Yopal, Barranquilla, San Andrés, Ibagué, Zipaquirá, Duitama, Montería, Pereira y Cali & 3.173 personas registradas & Identificar barreras, apropiación y soluciones territoriales basadas en IA & Diagnóstico regional participativo para política pública de IA & \parencite{minciencias_dialogos_ia_pdf_2024} \\
Política de IA & CONPES 4144 de IA & DNP y CONPES & 14-feb-2025--2030 & 479.000 millones de pesos reportados & Nacional & Ecosistema colombiano de IA & Implementar más de 100 acciones de política pública hasta 2030 & Gobernanza de IA, adopción, infraestructura, talento y mitigación de riesgos & \parencite{vanegas_conpes_4144_2025,conpes_4144_ia_pdf_2025} \\
Política de IA & Participación MinCiencias en CONPES 4144 & MinCiencias & 2025--2030 & Incluida en financiación CONPES & Nacional & Ecosistema de I+D+i, entidades públicas y sociedad & Ejecutar acciones asignadas a MinCiencias dentro de la política nacional de IA & Consejo asesor de expertos, centros de datos a hiperescala, transferencia de tecnologías basadas en IA, beneficios tributarios y apropiación social & \parencite{conpes_4144_ia_pdf_2025} \\
IA pública & Guía Ética de IA pública & Presidencia, DNP, MinTIC, MinCiencias y CRC & 31-dic-2025 / 07-ene-2026 & No especificado & Entidades públicas nacionales y territoriales & Servidores públicos, entidades y ciudadanía usuaria de sistemas públicos & Orientar adopción responsable de IA en el sector público & Principios, orientaciones y lineamientos para implementación, desarrollo y uso de IA pública & \parencite{dnp_guia_etica_ia_2025,mintic_guia_etica_ia_2026} \\
IA y cuántica & ColombIA Inteligente & MinCiencias & 2025--2026 & 24.000 millones de pesos para convocatoria 2026 & Nacional, con vocación territorial & Universidades, centros de I+D+i, alianzas regionales y talento de posgrado & Financiar soluciones de IA y tecnologías cuánticas para retos sociales, productivos y ambientales & Hasta 2.000 millones de pesos por proyecto; IA y cuánticas; duración máxima de 14 meses; articulación con política nacional de IA & \parencite{minciencias_rendicion_2024_2025_pdf} \\
IA y cuántica & ColombIA Inteligente, Convocatoria 966 & MinCiencias & 25-abr-2025 / 18-jun-2025 & 8.377 millones de pesos programados & Nacional & Alianzas de investigación, desarrollo tecnológico e innovación & Financiar investigación aplicada en IA y tecnologías cuánticas & Apertura 25-abr-2025; cierre 18-jun-2025; resultados previstos para septiembre de 2025; tecnologías cuánticas e IA & \parencite{minciencias_rendicion_2024_2025_pdf} \\
IA y ciencias del espacio & ColombIA Inteligente, Convocatoria 950 & MinCiencias & 2024 & 6.627 millones de pesos de MinCiencias y 1.889 millones de contrapartida & Nacional, con enfoque en problemas territoriales & 5 proyectos de I+D+i; 40 jóvenes investigadores; 7 estudiantes de maestría; 3 estancias posdoctorales & Financiar soluciones con IA y ciencias del espacio para problemas territoriales & Investigación aplicada, desarrollo tecnológico, innovación, IA y ciencias del espacio & \parencite{minciencias_rendicion_2024_2025_pdf} \\
Robótica educativa & Colombia Robótica, campamentos STEAM & MinCiencias & 2024 & 806,2 millones de pesos & Antioquia, Cundinamarca y Tumaco & 1.200 niños, niñas y adolescentes y 80 docentes & Promover apropiación temprana de tecnologías digitales y vocaciones CTeI & Programación, robótica, diseño digital, realidad virtual e inteligencia artificial & \parencite{minciencias_rendicion_2024_2025_pdf} \\
Robótica educativa & Colombia Robótica, laboratorios STEAM & MinCiencias & 2025 & 4.800 millones de pesos & Tumaco & Instituciones educativas y estudiantes beneficiarios de laboratorios & Adecuar espacios de aprendizaje STEAM y vocaciones CTeI & 15 laboratorios STEAM; equipamiento tecnológico; mobiliario especializado; kits de robótica y material didáctico & \parencite{minciencias_rendicion_2024_2025_pdf} \\
Ciencia para la paz & Catatumbo, Ciencia para la Paz & MinCiencias & 26-ene-2025 a 07-feb-2025 & 133,220 millones de pesos & Catatumbo & 430 niños, niñas, adolescentes y jóvenes & Responder a la emergencia territorial con apropiación científica y tecnológica & 500 kits de robótica y electrónica maker; 26 talleres de robótica, programación, electrónica maker, astronomía y cartografía participativa & \parencite{minciencias_rendicion_2024_2025_pdf} \\
Mujeres en ciencia & Orquídeas, mujeres en la ciencia & MinCiencias & 2023--2024 & 20.611 millones de pesos en 2023 y 28.702 millones de pesos en 2024 & Nacional, con enfoque territorial y de género & Doctoras y jóvenes investigadoras; 107 proyectos en 2023 y meta de 119 proyectos en 2024 & Vincular mujeres investigadoras a proyectos de I+D+i & Proyectos de investigación, formación de alto nivel, enfoque territorial, bioeconomía, IA y ciencia para la paz & \parencite{minciencias_informe_gestion_2024_pdf} \\
Investigación fundamental & Investigación Fundamental, Convocatoria 937 & MinCiencias & 2023--2024 & 19.032 millones de pesos & Nacional & 24 proyectos de I+D+i seleccionados & Fortalecer capacidades de investigación fundamental y producción científica & Proyectos con aprendizaje de máquina e IA aplicada en materiales y salud & \parencite{minciencias_informe_gestion_2024_pdf} \\
Seguridad digital & CiberPaz & MinTIC y Universidad de Pamplona en formaciones & 2025 & 8.000 millones de pesos para sensibilizaciones y 6.347 millones de pesos para formaciones & Nacional & Ciudadanía general; mujeres; población LGBTIQ+; ROM; indígenas; NARP; rurales; personas mayores & Promover uso seguro, ético y productivo de TIC y formar en IA aplicada & Talleres, festivales, hackathones, alfabetización mediática, fake news, cursos por WhatsApp y módulo virtual con enfoque en IA & \parencite{mintic_informe_gestion_molina_2025} \\
Seguridad digital & Estrategia Nacional de Seguridad Digital & MinTIC y ColCERT & 24-jun-2025 & No especificado & Nacional & Entidades públicas, sector privado y ciudadanía & Responder a amenazas digitales crecientes y fortalecer resiliencia & Estrategia 2025--2027; ciberseguridad, gestión de riesgos, coordinación y resiliencia institucional & \parencite{mintic_seguridad_digital_2025} \\
Seguridad digital & Actualización del MSPI & MinTIC & 03-jun-2025 & No especificado & Entidades públicas & Entidades obligadas por lineamientos de seguridad y privacidad & Actualizar el Modelo de Seguridad y Privacidad de la Información & Resolución 02277; alineación con ISO/IEC 27001:2022 & \parencite{mintic_mspi_resolucion_02277_2025} \\
Defensa digital & Arquitectura de ciberdefensa sector defensa & Ministerio de Defensa, Fuerzas Militares y Policía Nacional & 2025 & 9.006 millones de pesos en herramientas de seguridad; 4.142 millones de pesos en monitoreo y forense digital & Nacional e infraestructuras críticas del sector defensa & Entidades del sector defensa, sistemas críticos y ciudadanía protegida & Pasar de un esquema reactivo a capacidades anticipativas de monitoreo, prevención y respuesta & Sistemas SIEM, WAF, firewalls dinámicos, segmentación de tráfico, Anti-DDoS y capacidades de monitoreo y forense digital & \parencite{mindefensa_audiencia_rendicion_2025,cogfm_informe_ejecutivo_2025} \\
Defensa digital & Tecnologías de decepción cibernética & Ministerio de Defensa & 2025 & 2.369 millones de pesos & Infraestructura digital del sector defensa & Equipos de ciberdefensa y entidades protegidas & Anticipar y caracterizar comportamientos maliciosos antes de que afecten sistemas reales & Entornos simulados, identificación de patrones de ataque y fortalecimiento de respuesta temprana & \parencite{mindefensa_audiencia_rendicion_2025} \\
Defensa digital / democracia & PMU Ciber electoral & Ministerio de Defensa y entidades participantes & 2025 & No especificado & Nacional, procesos electorales & Organización electoral, entidades de seguridad y ciudadanía & Proteger procesos democráticos frente a riesgos digitales & Puesto de Mando Unificado Ciber con participación de ocho entidades; coordinación interinstitucional para gestión de riesgo electoral digital & \parencite{mindefensa_audiencia_rendicion_2025,mintic_seguridad_digital_2025} \\
Defensa digital / Policía & Centro Cibernético Policial y CAI Virtual & Policía Nacional de Colombia & 2025 & No especificado & Nacional y ciberespacio & Ciudadanía, entidades públicas y privadas, niñas, niños y adolescentes & Ejercer autoridad efectiva en el ciberespacio, atender incidentes y prevenir delitos informáticos & 17.321 bloqueos preventivos; 11.373 incidentes atendidos por CAI Virtual; 245 alertas preventivas; 67 riesgos identificados; 2.149 intentos de phishing detectados & \parencite{policia_rendicion_logros_retos_2025} \\
Defensa digital / Policía & Estrategia ESCIB--C2 & Policía Nacional de Colombia & 2025 & No especificado & Nacional & Víctimas de delitos informáticos y explotación sexual infantil & Judicializar delitos informáticos y explotación sexual infantil en entornos digitales & 46 capturas reportadas: 31 por delitos informáticos y 15 por explotación sexual infantil & \parencite{policia_rendicion_logros_retos_2025} \\
Defensa digital / niñez & Plan Pescador & Policía Nacional de Colombia & No especificado & No especificado & Nacional y cooperación internacional & Niñas, niños y adolescentes en entornos digitales & Proteger derechos de la niñez frente a explotación infantil y delitos en línea & Incautación de dispositivos, análisis forense digital, judicialización de redes criminales y cooperación con HSI y NCMEC & \parencite{policia_plan_pescador_ninez_digital,policia_rendicion_logros_retos_2025} \\
Soberanía tecnológica defensa & Software propio del Ejército y COGFM & Ejército Nacional y Comando General de las Fuerzas Militares & 2025 & Ahorros reportados: 4.200 millones de pesos en CAOCC PDF; 256 millones en ORIÓN; 384 millones en SIJEN; 108 millones en SISPRO; 500 millones en SICOE 2.0; 960 millones en casinos & Sector defensa & Unidades militares, áreas administrativas y tomadores de decisión & Reducir dependencia de proveedores externos y alojar información estratégica en servidores institucionales & CAOCC PDF, ORIÓN, SIJEN, modernización de SISPRO, SICOE 2.0 para analítica en tiempo real y sistemas de gestión de casinos & \parencite{mindefensa_audiencia_rendicion_2025,cogfm_informe_ejecutivo_2025} \\
Soberanía tecnológica defensa & Reclutamiento interoperable y SGDEA & Fuerzas Militares / sector defensa & 2025 & No especificado & Nacional e institucional & Ciudadanía usuaria del sistema de reclutamiento y entidades del sector defensa & Modernizar trámites, gestión documental e interoperabilidad institucional & Sistema de reclutamiento interoperable con Registraduría Nacional y plataforma SGDEA para gestión documental interoperable entre fuerzas & \parencite{mindefensa_audiencia_rendicion_2025,cogfm_informe_ejecutivo_2025} \\
Soberanía tecnológica defensa & PadillApp y chatbots institucionales & Armada Nacional y Fuerza Aeroespacial Colombiana & 2023--2025 & No especificado & Institucional & Personal de la Armada, Fuerza Aeroespacial y usuarios de servicios institucionales & Mejorar interacción digital interna y atención institucional mediante herramientas propias & Aplicación móvil PadillApp en la Armada Nacional y chatbots institucionales en la Fuerza Aeroespacial Colombiana & \parencite{arc_informe_rendicion_2023,fac_informe_rendicion_2023} \\
Infraestructura defensa & Redes, respaldo de datos y equipos sector defensa & Ministerio de Defensa, Fuerzas Militares y COGFM & 2025 & 8.598 millones de pesos en redes integradas; 2.734 millones en redes de alta velocidad; 2.869 millones en FREQUENTIS & Infraestructura tecnológica del sector defensa & Unidades operativas y administrativas del sector defensa & Fortalecer continuidad operativa, reducir latencia y robustecer la base material de soberanía digital & Redes de comunicación hasta 400 Gbps; respaldo de datos de 40 TB a 223 TB; sistemas FREQUENTIS; distribución de 1.849 computadores & \parencite{mindefensa_audiencia_rendicion_2025,cogfm_informe_ejecutivo_2025} \\
Infraestructura defensa marítima & Comunicaciones y vigilancia marítima digital & Armada Nacional de Colombia & 2023--2025 & 3.700 millones de pesos en vigilancia marítima; 1.498 millones de pesos en videoconferencia institucional & Zonas marítimas y unidades navales, incluidas zonas remotas & Armada Nacional, seguridad marítima y comunidades atendidas por presencia naval & Fortalecer comunicaciones, vigilancia marítima y presencia digital en zonas remotas & Comunicaciones satelitales Iridium, fibra óptica en zonas remotas, modernización de vigilancia marítima y sistemas de videoconferencia institucional & \parencite{arc_informe_rendicion_2023,mindefensa_audiencia_rendicion_2025} \\
Espacio y defensa & Satélite FACSAT-2 & Fuerza Aeroespacial Colombiana & 2023--2025 & No especificado & Nacional, con aplicaciones de inteligencia y gestión del riesgo & Sector defensa, gestión del riesgo y tomadores de decisión & Producir información geoespacial propia para inteligencia y gestión del riesgo & 725 imágenes capturadas y 659 procesadas para misiones de inteligencia y gestión del riesgo & \parencite{fac_informe_rendicion_2023,mindefensa_audiencia_rendicion_2025} \\
Aeronaves autónomas & Proyecto Ala Zagi & Fuerza Aeroespacial Colombiana / sector defensa & 2023--2025 & No especificado & Nacional e institucional & Unidades de vigilancia y reconocimiento & Desarrollar aeronaves autónomas de bajo costo para vigilancia & Aeronaves autónomas de bajo costo, vigilancia y apoyo a capacidades de reconocimiento & \parencite{fac_informe_rendicion_2023,mindefensa_audiencia_rendicion_2025} \\
Drones defensa & Flota de drones de inteligencia, reconocimiento y operaciones & Ministerio de Defensa y Fuerzas Militares & 2025--2026 & 30.000 millones de pesos gestionados para tecnologías anti-drones; 126.200 millones priorizados para capacidades frente a amenazas aéreas no tripuladas & Nacional y zonas estratégicas de operación & Fuerza Pública, seguridad territorial y población protegida & Incrementar inteligencia, reconocimiento, operación y respuesta frente a amenazas no tripuladas & Flota de 966 drones; 66 modelos de alta capacidad; incremento reportado de 147\% en poder de fuego institucional & \parencite{mindefensa_audiencia_rendicion_2025,mindefensa_100_logros_2026} \\
Antidrones defensa & Sistema NESHER V4 & Ministerio de Defensa / Fuerza Pública & 2025--2026 & Incluido en capacidades anti-drones y defensa frente a amenazas no tripuladas & Zonas estratégicas & Fuerza Pública e infraestructuras protegidas & Desarrollar capacidades propias de detección e inhibición contra amenazas aéreas no tripuladas & Sistema NESHER V4 con patentes propias; 34 sistemas ensamblados; 35 detectores desplegados; neutralización de amenazas en campo & \parencite{mindefensa_audiencia_rendicion_2025,mindefensa_100_logros_2026} \\
Gestión pública defensa & SIDEH, SEGET y SIIPO & Ministerio de Defensa Nacional & 2025 & No especificado & Nacional y territorios priorizados por seguridad, derechos humanos y posconflicto & Entidades del sector defensa, territorios priorizados y población beneficiaria de rutas de protección & Fortalecer gestión territorial, seguimiento a derechos humanos y compromisos del posconflicto con sistemas de información & SIDEH con seguimiento a 1.399 activaciones de rutas de protección; SEGET para priorización de inversiones; SIIPO para seguimiento de compromisos de implementación & \parencite{mindefensa_siipo_implementacion_2025,mindefensa_audiencia_rendicion_2025} \\
CTeI defensa & Política sectorial de CTeI, Resolución 2443 de 2024 & Ministerio de Defensa y Fuerza Pública & 2024--2025 & No especificado & Sector defensa & Fuerza Pública, ecosistema de innovación y capacidades tecnológicas nacionales & Cerrar brechas tecnológicas con enfoque de misiones y fortalecer propiedad intelectual y transferencia tecnológica & Gestión de 94 procesos de propiedad intelectual; 47 registros obtenidos; transferencia tecnológica; mantenimiento aeronáutico; comunicaciones; entrenamiento táctico avanzado & \parencite{mindefensa_audiencia_rendicion_2025,cogfm_informe_ejecutivo_2025} \\
Tecnologías emergentes defensa & Computación cuántica y criptografía post-cuántica & Ministerio de Defensa / Fuerza Pública & 2025 & No especificado & Sector defensa & Capacidades estratégicas del Estado & Explorar tecnologías de frontera para seguridad, defensa y resiliencia futura & Exploración de computación cuántica y criptografía post-cuántica articulada a la política sectorial de CTeI & \parencite{mindefensa_audiencia_rendicion_2025,cogfm_informe_ejecutivo_2025} \\
Defensa marítima y ciencia & ARC Simón Bolívar y capacidades científicas marítimas & Armada Nacional de Colombia y Ministerio de Defensa & 2023--2025 & No especificado & Dominio marítimo colombiano & Armada Nacional, comunidad científica, seguridad marítima y gestión ambiental & Extender la soberanía digital al dominio marítimo y científico mediante generación de conocimiento propio & Buque ARC Simón Bolívar, simulación inmersiva, sistemas de gestión de combate de fabricación nacional, robótica marina, modelos predictivos SAR y monitoreo ambiental/marítimo & \parencite{arc_informe_rendicion_2023,mindefensa_audiencia_rendicion_2025} \\
Servicios digitales defensa & SALUD.SIS, PQR virtuales y plataformas educativas & Fuerza Aeroespacial Colombiana y sector defensa & 2023--2025 & 9.000 millones de pesos en SALUD.SIS & Sector defensa & Personal del sector defensa, usuarios de salud institucional y ciudadanía usuaria de canales digitales & Modernizar servicios institucionales, salud, formación y rendición de cuentas mediante canales digitales & 98,5\% de peticiones FAC por canales virtuales; más de 575.000 visitas al portal; plataformas educativas con más de 300.000 usuarios; telemedicina y telemonitoreo; bibliotecas digitales para más de 60 unidades & \parencite{fac_informe_rendicion_2023,mindefensa_audiencia_rendicion_2025} \\
Regulación de conectividad & Internet comunitario fijo & Presidencia y MinTIC & 30-jun-2023 & No especificado & Nacional & Comunidades organizadas y usuarios de Internet comunitario fijo & Definir condiciones para prestación del servicio de Internet comunitario fijo & Decreto 1079; habilitación normativa para prestación comunitaria de conectividad & \parencite{funcionpublica_decreto_1079_2023} \\
Regulación de conectividad & Regulación diferencial ICF & CRC & 28-mar-2025 & No especificado & Nacional & Proveedores y usuarios de Internet Comunitario Fijo & Crear medidas diferenciales para provisión de Internet Comunitario Fijo & Resolución CRC 7712; reglas diferenciales para prestación comunitaria & \parencite{crc_resolucion_7712_2025} \\
Estado digital & CMIC & MinTIC & 20-dic-2024 & No especificado & Nacional & Usuarios y prestadores de servicios de comunicaciones & Monitorear e inspeccionar servicios de comunicaciones & Centro de Monitoreo e Inspección de Servicios de Comunicaciones; reporte de 1.378 estaciones base 5G al tercer trimestre de 2024 & \parencite{mintic_cmic_telecomunicaciones_2024} \\
Estado digital & Rendición MinTIC 2024 & MinTIC & 09-dic-2024 & No especificado & Nacional & Ciudadanía y usuarios de política TIC & Reportar avances anuales de conectividad, infraestructura y apropiación digital & Cerca de tres millones de nuevas personas conectadas; 1.353 estaciones base 5G; 909 Juntas de Internet; 78.060 computadores; 2.375 antenas & \parencite{mintic_rendicion_cuentas_2024} \\
Estado digital & Servicios Ciudadanos Digitales & MinTIC y Agencia Nacional Digital & 2025 & 19.000 millones de pesos & Nacional & Ciudadanía y entidades públicas & Sostener y evolucionar el ecosistema de interacción digital Estado-ciudadano & Soporte, evolución, mantenimiento, operación, uso y apropiación de plataformas SCD bajo Convenio 1225-2025 & \parencite{mintic_informe_gestion_molina_2025} \\
Estado digital & Mi Colombia Digital & MinTIC & 2022--2025 & 16.195 millones de pesos en Gobierno del Cambio; 10.000 millones de pesos en 2025 & 1.103 municipios y entidades territoriales & Entidades públicas territoriales y ciudadanía & Proveer sedes electrónicas y herramientas colaborativas para gobierno digital & 9.774 sedes electrónicas activas; 35.903 correos entregados; integración con trámites, SECOP, SIGEP, carpeta ciudadana, autenticación digital y datos & \parencite{mintic_informe_gestion_molina_2025} \\
Estado digital & Carpeta Ciudadana y REDAM & Agencia Nacional Digital & 2024 & No especificado & Nacional & Ciudadanía usuaria de servicios digitales del Estado & Fortalecer servicios digitales ciudadanos y registros administrativos & 5.339.021 certificados expedidos, 1.738 deudores registrados, mejoras operativas y capacitaciones & \parencite{and_informe_gestion_2024} \\
Estado digital tributario & Servicios digitales DIAN & DIAN & 30-ene-2026 & No especificado & Nacional & 35.256.814 personas con información exógena; 17.365.616 con reportes de facturas & Mejorar servicios tributarios digitales e interacción con contribuyentes & Información exógena, validación de identidad, reportes de facturas y servicios sugeridos & \parencite{dian_informe_gestion_2025} \\
Estado digital tributario & Sistema de Facturación Electrónica & DIAN & 31-ene-2026 & No especificado & Nacional & Facturadores electrónicos, usuarios de RADIAN y documentos equivalentes & Consolidar sistemas digitales de facturación y trazabilidad tributaria & Series y cortes de habilitados en factura electrónica, RADIAN y documentos equivalentes & \parencite{dian_cifras_sfe_2026} \\
Datos públicos & Ley 2335 de estadísticas oficiales & Congreso y DANE & 03-oct-2023 & No especificado & Nacional & Sistema Estadístico Nacional, entidades públicas y ciudadanía & Fortalecer marco de operaciones estadísticas, registros administrativos y datos oficiales & Ley 2335; reglas para operaciones estadísticas, registros administrativos y datos del Sistema Estadístico Nacional & \parencite{funcionpublica_ley_2335_2023} \\
Protección de datos e intimidad & Ley 2300 de protección de intimidad & Congreso & 10-jul-2023 & No especificado & Nacional & Consumidores y personas contactadas por cobranza o relaciones comerciales & Proteger intimidad en canales, horarios y frecuencia de contacto & Reglas de canales, horarios y periodicidad de contacto en cobranza y relaciones de consumo & \parencite{funcionpublica_ley_2300_2023} \\
Data centers & Campus Bogotá de Ascenty & Ascenty & 2023--2025 & No especificado & Bogotá-Región & Clientes empresariales, cloud y tecnología & Ampliar infraestructura de colocation y conectividad de grado industrial & 2 data centers en Colombia; 20 MW; 18.500 m2; carrier neutral; conectividad y colocation & \parencite{ascenty_colombia_data_centers} \\
Data centers & BG-02 y BG-03 de Aligned/ODATA & Aligned Data Centers & 2025--2026 & No especificado & Colombia / Bogotá-Región & Demanda enterprise, hyperscale y cargas de IA & Expandir capacidad de data center para alta densidad, cloud, IA y HPC & BG-01 existente; nuevas instalaciones BG-02 y BG-03 con 144 MW combinados; certificaciones y conectividad & \parencite{aligned_colombia_data_centers} \\
\end{longtable}
\endgroup
